\documentclass[final,5p,times,twocolumn]{elsarticle}

\usepackage[T1]{fontenc}
%\usepackage[utf8]{inputenc}
\usepackage{bm}
\usepackage{amsfonts}
\usepackage{graphicx}
\usepackage{slashed}
\usepackage{color}
\usepackage{cancel}
\usepackage{hyperref}
\usepackage{amsthm}
\usepackage{amssymb}
\usepackage{amsmath}
\usepackage{xparse}
\usepackage{xspace}
\usepackage{mathtools}
\usepackage{etextools}
\usepackage{ifthen}
\usepackage{bm}
\usepackage{empheq}
\usepackage[most]{tcolorbox}
\newcommand{\Tends}{\,\rightarrow\, }
\newcommand{\COM}[2]{\ensuremath{\left[{#1} ,{#2}\right]}\xspace}
\newcommand{\Order}[1]{\mathcal{O}\left(#1 \right)}
\DeclarePairedDelimiter{\MMbBraces}{\{\!\{\!\{}{\}\!\}\!\}}

\DeclarePairedDelimiter{\MbBraces}{\{\!\{}{\}\!\}}
\newcommand{\PB}[2]{\left\{#1,#2\right\}}
\newcommand{\QPB}[2]{\left\{#1,#2\right\}}
\newcommand{\MB}[2]{\MbBraces{#1,#2}}
%\newcommand{\MMB}[2]{\MMbBraces{#1,#2}}
\newcommand{\MMB}[2]{\PB{#1}{#2}_\mathcal{W}}
\newcommand{\ud}{\mathrm{d}}
\newcommand{\dW}{\cal{W}}
\newcommand{\uHam}{\mathbf{\Ham}}
\newcommand{\uW}{\mathbf{W}}
\newcommand{\uF}{\mathbf{F}}
\newcommand{\uG}{\mathbf{G}}
\newcommand{\uC}{\mathbf{C}}
\newcommand{\uK}{\mathbf{K}}
\newcommand{\uH}{\mathbf{H}}
\newcommand{\uO}{\mathbf{O}}
\newcommand{\uP}{\mathbf{P}}
\newcommand{\HPD}[3]{\ensuremath{\frac{\partial^{#3}#1}{{\partial#2}^{#3}}}\xspace}
\NewDocumentCommand{\PD}{o m m}{%
  \IfNoValueTF{#1}
    {\ensuremath{\frac{\partial #2}{\partial #3}}}%
    {\ensuremath{\frac{\partial^{#1} #2}{\partial #3^{#1}}}}%
  \xspace
}
\NewDocumentCommand{\TD}{o m m}{%
  \IfNoValueTF{#1}
    {\ensuremath{\frac{\ud #2}{\ud #3}}}%
    {\ensuremath{\frac{\ud^{#1} #2}{\ud #3^{#1}}}}%
  \xspace
}
\newcommand{\ui}{\mathrm{i}}
\newcommand{\FIX}[1]{\textcolor{red}{#1}}
\newcommand{\Gam}[1]
{{\sqrt{1 - \dfrac{{#1}^2}{c^2}}}}
\newcommand{\Lgn}{\mathcal{L}}
\newcommand{\Ham}{{H}}

\usepackage{amsmath,amssymb,bm}
\hypersetup{hidelinks}
\journal{Physics Letters A}

\begin{document}
%\enlargethispage{300pt}

\begin{frontmatter}

\title{Spinor Structure from Relativistic Mass-Shell Factorisation in Phase Space}

\author[1]{Mark Everitt\corref{cor1}}
\cortext[cor1]{Corresponding author}
\ead{m.j.everitt@physics.org}
\address[1]{Quantalytics, Spinney Hill Drive, Loughborough, LE11 3LB, UK}

\begin{abstract}
Spin is usually regarded as one of the most intrinsically quantum phenomena, while its lack of a natural analogue in classical physics presents an obstacle to phase-space deformation-quantisation accounts of its origin. 
We show that this difficulty may arise from imposing an overly restrictive scalar Hamiltonian structure on relativistic phase space. 
Requiring the complete massive mass-shell constraint to be represented by a single finite-dimensional expression that is linear in all four components of momentum forces its coefficient matrices to satisfy a Clifford algebra. 
The minimal complex representation of this algebra is four-dimensional. 
Requiring statistical completeness within the rank-two subspace selected by the mass-shell factor then leads to a four-by-four matrix-valued ensemble distribution.
At each on-shell momentum, the linear mass-shell operator selects a two-dimensional subspace, so a general classical ensemble is described by a $2 \times 2$ matrix before quantisation.
Projecting the Weyl-ordered Liouville equation into this subspace gives relativistic transport while preserving arbitrary populations and coherences and the remaining projector components determine the first deformation correction. 
Expansion of the matrix Moyal star commutator yields the symmetrised classical matrix Liouvillian at leading order.
If the stronger two-sided star constraints are imposed, they reproduce the left and right Dirac-Wigner equations and introduce $\hbar$ as the scale converting the dimensionless internal algebra into physical angular momentum.
These results suggest a non-quantum origin for spinor structure and a route to reconciling relativistic covariance, classical phase-space transport, and quantum spin.
\end{abstract}

\begin{keyword}
	Relativistic statistical mechanics \sep 
	covariant classical phase space    \sep 
	classical spinor structure 		   \sep 
	Clifford algebra                   \sep 
	mass-shell factorisation           \sep 
	deformation quantisation           \sep 
	Dirac--Wigner equations
\end{keyword}

\end{frontmatter}
Spin-$1/2$ is normally introduced into relativistic quantum mechanics through a spinor state space. 
We ask whether the underlying spinor structure can instead follow from compatibility between relativistic kinematics and statistical phase-space dynamics. 
This study was motivated by the view that any physical theory should allow statistical descriptions of ensembles without selecting a preferred observer.
This makes the covariant form of Liouville's equation a natural starting point, although its phase-space generator need not have the interpretation
of a coordinate-time energy~\cite{hakim2011introduction}. 
Canonical interacting particle models are also restricted by the Currie--Jordan--Sudarshan no-interaction theorem, which shows that, for a finite number of classical point particles with canonical coordinates whose positions transform covariantly, there is no non-trivial Hamiltonian interaction compatible with the full Poincar\'e symmetry (the fundamental symmetries of special relativity)~\cite{RevModPhys.35.350}. 

Phase-space formulations hold  the promise that quantum mechanics naturally emerges as a deformation quantisation of classical physics, replacing the Poisson bracket with the Moyal bracket and conventional distribution multiplication with the star product~\cite{GROENEWOLD1946405,Moyal_1949}. 
However, a persistent gap in this approach has been the difficulty of naturally introducing spin, since quantum mechanical spin-$1 \over 2$ has no direct analogue in classical mechanics. 
Previous work has introduced spin phase-space representations through $\mathrm{SU}(N)$ Wigner functions formulated in terms of a displaced-parity operator~\cite{PhysRevA.99.012115}.
More recently, Ji and Piasecki mapped the relativistic Wigner operator, through its Clifford decomposition, to a Koopman phase-space spinor with a classical Liouville limit~\cite{sym18061018}.
The present construction addresses the complementary question of how the Clifford algebra and its rank-two spinor structure follow directly from first-order mass-shell factorisation, before quantisation.
The argument developed below revisits Dirac's original derivation, but reformulates it in a classical phase-space language: 
the relativistic Liouville equation replaces Schr\"odinger's equation as the dynamical starting point, and spin degrees of freedom are encoded via a Clifford algebra-valued distribution. 
For a relativistic extension of Schr\"odinger's equation, Dirac argued that 
(i) it should make use of the mass shell, 
(ii) the dynamical equation, like the Schr\"odinger equation, should be first-order in time but also, inspired by issues with the Klein-Gordon equation, that the Hamiltonian should be first-order in space, and 
(iii) the four-momentum components should be replaced with their operator counterparts~\cite{Dirac28}. 
Unlike Dirac's argument, we will use the Liouville equation as the basis for the discussion, but apart from this, and the fact that no operator substitution is required, we impose the following three requirements: 
(a) the dynamics should be first order in spacetime and momentum derivatives 
(b) both sheets of the relativistic mass shell should be retained, and 
(c) the mass-shell constraint should be represented linearly. 
These requirements lead naturally to a Clifford factorisation. 
In $3+1$ dimensions its minimal complex representation is four-dimensional, while its on-shell eigenspaces have rank two.
In this sense, the algebra of spin arises kinematically from the compatibility of relativistic mass-shell factorisation with a statistical phase-space description.
At each on-shell momentum the factor selects a two-dimensional subspace. 
A statistical ensemble that retains all information within this subspace is therefore naturally matrix valued before any attempt at quantisation.
Deformation quantisation subsequently introduces the nonlocal star product and the scale $\hbar$. 
In this sense, factorisation supplies the dimensionless spinor algebra, whereas $\hbar$ supplies the physical scale of the angular momentum associated with spin.


We work on the extended phase space formed by the spacetime coordinates
$  
    x\,^\mu=(ct,\bm{x})
$ where Greek indices run from $0$ to $3$, repeated indices are summed, and we use the metric $\eta_{\mu\nu}=\operatorname{diag}(1,-1,-1,-1)$. 
Proper time $\tau$ is defined along any given worldline by $\ud s = c\,\ud\tau $ where the line element $\ud s^2 = \eta_{\mu\nu} \, \ud x\,^\mu \ud x\,^\nu = c^2 \ud t^2 - \ud\Vec{x}\,^2
$. 
The three-velocity is $v^i \equiv \ud x\,^i/\ud t$, and $\Vec{v}$ denotes the corresponding three-vector. 
The physical three-momentum is $\Vec{p}=\gamma m\Vec{v}$, where
$\gamma=(1-\Vec{v}\,^2/c^2)^{-1/2}$, and the corresponding
future-directed four-momentum is $p\,^\mu_+=(E_p/c,\Vec{p})$, with
$E_p=\sqrt{c^2\Vec{p}\,^2+m^2c^4}$. To retain the complete mass shell,
we extend the phase-space description to
$p^\mu_\sigma=(\sigma E_p/c,\Vec{p})$, where the mass-shell label $\sigma=\pm1$. We refer to the
resulting eight-dimensional space of four-position and four-momentum as phase spacetime.



The relativistic Liouville equation is defined in terms of $H(x,p)$, a scalar covariant phase-space generator (rather than necessarily a coordinate-time energy) and $W(x,p)$, a scalar ensemble distribution.  Collisionless evolution is expressed through the vanishing extended Poisson
bracket~\cite{MARSDEN198629,hakim2011introduction}
\begin{equation}
  \mathcal{L}_H(W)
  \equiv
  \PB{H}{W}
  \equiv 
  \PD{H}{p_\nu}
  \PD{W}{x^\nu}
  -
  \PD{W}{p_\nu}
  \PD{H}{x^\nu}
  =0.
  \label{eq:relL}
\end{equation}
Because this equation is formulated covariantly on extended phase space, we do not identify $H$ with the coordinate-time energy
$\gamma mc^2$. Instead, we begin from the invariant mass-shell
constraint
\begin{equation}
  p\,^\mu p_\mu=m^2c^2,
  \label{eq:mass-shell}
\end{equation}
Solving the constraint only for $p^0=E_{\bm p}/c$ would select the
future-directed sheet. We instead seek a single Lorentz-covariant
factor that is linear in all four components of $p^\mu$. Let
$$
    K(p)=A^\mu p_\mu
$$
be such a finite-dimensional matrix-valued expression. For it to
represent the complete quadratic mass-shell constraint, its square
must satisfy
\begin{equation}
    K(p)^2
    =
    A^\mu A^\nu p_\mu p_\nu
    =
    \left(p^\mu p_\mu\right){1_N}
    \qquad
     \forall p .
\end{equation}
Because $p_\mu p_\nu$ is symmetric, only the anticommutator of the
coefficient matrices contributes:
\begin{equation}
    K(p)^2
    =
    \frac{1}{2}
    \left\{A^\mu,A^\nu\right\}p_\mu p_\nu .
\end{equation}
Comparison of the coefficients of $p_\mu p_\nu$ therefore requires
\begin{equation}
    \left\{A^\mu,A^\nu\right\}
    =
    2\eta^{\mu\nu}{1_N}.
\end{equation}
Thus the coefficient matrices necessarily generate the Clifford
algebra. Writing $A^\mu\equiv\gamma^\mu$, its minimal complex
representation in $3+1$ dimensions is four-dimensional, and the
massive constraint factorises as
\begin{equation}
    \left(\gamma^\mu p_\mu-mc {1}_4\right)
    \Big(\gamma^\nu p_\nu+mc {1}_4\Big)
    =
    \left(p^\mu p_\mu-m^2c^2\right){1}_4 .  \label{eq:factorisation}
\end{equation}
On the mass shell, the product in Eq.~(6) vanishes, and each linear factor has a nontrivial kernel. Importantly, the components of \(p^\mu\) remain commuting classical variables where no operator substitution has been made. 
The construction mirrors Dirac's algebraic factorisation, but it is carried out entirely on relativistic phase space. 
We therefore introduce the matrix-valued covariant phase-space generator
\begin{equation}
  \uH_{\mathrm{free}}(p)
  =c(\gamma\,^\mu p_\mu-mc 1_4).
  \label{eq:matrix-generator}
\end{equation}
The linear factor acts in the minimal four-dimensional representation of the Clifford algebra. 
To formulate a statistical theory on this representation, we introduce a $4\times4$ matrix-valued covariant phase-space function $\uW(x,p)$, which we call the matrix ensemble distribution. 
The need for a general matrix distribution does not follow from the representation alone, rather it follows below when we require the statistical description to retain all information within the subspace selected by the factor constraint.
One extension of Eq.~\eqref{eq:relL} to matrix form is
\begin{equation}
  \mathcal{L}_{\uH}(\uW)
  \equiv
  \PD{\uH}{p_\nu}
  \PD{\uW}{x^\nu}
  -
  \PD{\uW}{p_\nu}
  \PD{\uH}{x^\nu}
  \label{MLE}
\end{equation}
which we note retains the ordering of the usual, but not necessary, expression of the Poisson bracket.
Under a Lorentz transformation $\Lambda$, the distribution transforms
as
\begin{equation}
  \uW(x,p)\longmapsto
  S(\Lambda)
  \uW(\Lambda^{-1}x,\Lambda^{-1}p)
  S^{-1}(\Lambda).
\end{equation}
Because the factors are matrix valued, their order matters. The bracket in Eq.~\eqref{MLE} is bilinear and antisymmetric, but it does not in general satisfy the Jacobi identity or the Leibniz rule expected of a Lie bracket. 
It should therefore be regarded as one possible ordered extension of Eq.~\eqref{eq:relL} rather than as a Poisson bracket on a matrix algebra.

Because $\uH(x,p)$ and $\uW(x,p)$ no longer commute there are alternative brackets that would reproduce the normal Poisson bracket under the circumstances that $\uH$ and $\uW$ commute. 
For example consider the symmetric (Weyl-ordered) Liouvillian  
\begin{align}
    &
    \mathcal{L}^{\mathcal{W}}_\uH(\uW) \equiv
   \label{extendedbracket} 
   \\  & \ \ \ \ \ \ \ \ \ \ \ \ \ \ 
    \nonumber
    \frac{1}{2}
    \left[
    \left(
	\PD{\uH}{p_\nu}
	\PD{\uW}{x^\nu}
	+
	\PD{\uW}{x^\nu}
	\PD{\uH}{p_\nu}
	\right)
    -
	\left(
	\PD{\uH}{x^\nu}
	\PD{\uW}{p_\nu}
	+
	\PD{\uW}{p_\nu}
	\PD{\uH}{x^\nu}
	\right)
    \right]
%	  \right]
\end{align}
Algebraically, it is neither a Lie bracket nor a derivation. We will later find that it is nevertheless singled out by deformation quantisation, where we will see that it appears as the local $\Order{\hbar^0}$ term of the matrix star commutator. 
We do not yet impose the vanishing of the full matrix expression as the admissible classical matrix space should first be identified from the factor constraint.

Following~\cite{thaller1992dirac,BjorkenDrell1964RQM} we introduce the covariant mass-shell projectors
\begin{equation}
    \Pi_s(p)
    =
    \frac{1}{2}
    \left(
        1_4+s\frac{\gamma^\mu p_\mu}{mc}
    \right),
\end{equation}
Here $s=\pm 1$ labels the eigenspaces of $\gamma\,^\mu p_\mu/mc$ and should not be confused with the mass-shell label $\sigma=\operatorname{sgn}(p^0)$.
Note that~\cite{BjorkenDrell1964RQM} used the notation $\Lambda_\pm(p)$ which we have changed to avoid confusion with Lorentz transformations.
For any on-shell momentum $p_\sigma$, both $\Pi_+(p_\sigma)$ and $\Pi_-(p_\sigma)$ have rank
two. 
Moreover,
\begin{equation}
  \uH_{\mathrm{free}}(p_\sigma)\Pi_s(p_\sigma)
  =mc^2(s-1)\Pi_s(p_\sigma).
  \label{eq:projector-kernel}
\end{equation}
The zero-eigenvalue constraint therefore selects $s=+1$ on either
mass-shell sheet, while $\sigma$ distinguishes the two energy
branches.
These mass-shell projectors obey 
    $\Pi_s^2         = \Pi_s$, 
    $\Pi_s\Pi_r      = \delta_{sr}\Pi_s$,   
    and $\Pi_+ + \Pi_-   = 1_4$
when $p^\mu p_\mu=m^2c^2$.

At each on-shell momentum, the factor consequently selects the rank-two subspace
\begin{equation}
    \mathcal{E}_{p_\sigma}
    \equiv
    \ker \uH_{\mathrm{free}}(p_\sigma)
    =
    \operatorname{im}\Pi_+(p_\sigma).
    \label{eq:constraint-subspace}
\end{equation}
We now impose statistical completeness within $\mathcal{E}_{p_\sigma}$ so that the ensemble description must be capable of representing arbitrary populations and coherences within this rank-two subspace.
Accordingly, $W_0$ must allow an arbitrary $2\times2$ matrix $\uF_\sigma$ acting within $\mathcal{E}_{p_\sigma}$.
In the standard Dirac representation, the projectors at rest are
\begin{equation}
    \Pi_+(p_+)
    =
    \begin{pmatrix}
        1_2 & 0_2\\
        0_2 & 0_2
    \end{pmatrix},
    \qquad
    \Pi_+(p_-)
    =
    \begin{pmatrix}
        0_2 & 0_2\\
        0_2 & 1_2
    \end{pmatrix}.
\end{equation}
The corresponding distributions therefore have the forms
\begin{equation}
  \hspace{-4pt}  
  \uW_0(x,p_+)
    =
    \begin{pmatrix}
        \uF_+(x,\bm p) & 0_2\\
        0_2 & 0_2
    \end{pmatrix},
    \uW_0(x,p_-)
    =
    \begin{pmatrix}
        0_2 & 0_2\\
        0_2 & \uF_-(x,\bm p)
    \end{pmatrix}
    \label{eq:projected-classical-ansatz}
\end{equation}
where $\uF_\sigma$ is an arbitrary $2\times2$ matrix distribution acting within the selected subspace.
The position of the nonzero $2\times2$ block depends on the mass-shell sheet and the chosen basis. (note that for general nonzero momentum, however, the boost mixes the upper and lower components, so neither simple block form holds in the fixed Dirac basis).
Using $\uW_0=\Pi_{+}\uW_0\Pi_{+}$ together with $\uH_{\mathrm{free}}\Pi_{+} =\Pi_{+}\uH_{\mathrm{free}}=0_4$, we obtain
\begin{equation}
    \uH_{\mathrm{free}}\uW_0
    =0_4
    =\uW_0\uH_{\mathrm{free}} \text{\ \  and \ \  }
    \COM{\uH_{\mathrm{free}}}{\uW_0}=0_4.
    \label{eq:classical-commutator-constraint}
\end{equation}
This commutator is a classical compatibility condition between the distribution and the factor constraint and introduces no scale of action. If both eigenspaces of the factor are retained, the weaker condition is $\uW_0=\sum_s\Pi_s\uW_0\Pi_s$, which is equivalent on shell to $\COM{\uH_{\mathrm{free}}}{\uW_0}=0_4$. The zero-eigenvalue constraint selects the $\Pi_+$ term, while the mass-shell label $\sigma$ continues to distinguish the two energy sheets.

In the positive-energy rest frame where $p^\mu=(mc,\Vec{0})$, in the standard Dirac representation, we see
$$
        \Pi_+(mc,\mathbf{0})
        =
        \begin{pmatrix}
        1_2 & 0\\
        0 & 0
        \end{pmatrix}
        \text{\ \  and \ \ }
        \Pi_-(mc,\mathbf{0})
        =
        \begin{pmatrix}
        0 & 0\\
        0 & 1_2
        \end{pmatrix}.
$$
Crucially note that this rank-two structure also carries the dimensionless algebra of spatial rotations. 
Defining
\begin{equation}
  \Sigma^i
  =
  \frac{i}{4}\epsilon^{ijk}
  \COM{\gamma^j}{\gamma^k},
  \label{eq:dimensionless-spin}
\end{equation}
one finds, in the standard Dirac representation, $\Sigma^i=\operatorname{diag}(\sigma^i,\sigma^i)$. 
Thus, in the rest frame, the restriction of \(\Sigma^i\) to either rank-two sector furnishes the usual two-dimensional Pauli algebra. 
For general on-shell momentum \(p_\sigma\), the corresponding internal generators are obtained by transporting the rest-frame generators with a standard boost. 
The internal algebra is therefore preserved, although its embedding in the four-dimensional Clifford representation depends on momentum.
Importantly note that the compatibility condition does not reduce $\uF_\sigma$ to a scalar as arbitrary populations and coherences within this degenerate rank-two sector commute with $\uH_{\mathrm{free}}$. 
These internal coherences should not be confused with the off-diagonal blocks $\Pi_+\uW\Pi_-$ and $\Pi_-\uW\Pi_+$ that mix the two eigenspaces.
As in the Dirac theory, $\Sigma^i$ is dimensionless and \emph{no physical angular-momentum scale has yet been introduced}.
The following identities will be useful in simplifying projected equations
\begin{equation}
    \Pi_s \gamma^\mu \Pi_s = s \frac{p^\mu}{mc}\Pi_s \text{ and }\Pi_s \gamma^\mu \Pi_{-s}= \left(\gamma^\mu + s \frac{p^\mu}{mc} \right)\Pi_{-s}
\end{equation}
The factor constraint identifies which matrix components are allowed at leading classical order. Let
\begin{equation}
    \mathsf{D}_{\uH_{\mathrm{free}}}[\mathbf{X}]
    \equiv
    \sum_{s=\pm1}\Pi_s\mathbf{X}\Pi_s
    \label{eq:block-projection}
\end{equation}
denote the operation that retains only the blocks acting within each eigenspace of $\uH_{\mathrm{free}}$ of some arbitrary $\mathbf{X}$ (i.e. the $\mathsf{D}$iagonal blocks relative to the eigenspaces of $\uH_{\mathrm{free}}$). These are precisely the blocks that commute with $\uH_{\mathrm{free}}$. The constraint and transport equations can then be separated as
\begin{equation}
    \COM{\uH_{\mathrm{free}}}{\uW_0}=0_4,
    \qquad
    \mathsf{D}_{\uH_{\mathrm{free}}}\!\left[
        \mathcal{L}^{\mathcal{W}}_{\uH_{\mathrm{free}}}(\uW_0)
    \right]
    =0_4.
    \label{eq:projected-classical-system}
\end{equation}
For the zero-eigenvalue sector this reduces to
\begin{equation}
    \uH_{\mathrm{free}}\uW_0=0_4=\uW_0\uH_{\mathrm{free}},
    \qquad
    \Pi_+\mathcal{L}^{\mathcal{W}}_{\uH_{\mathrm{free}}}(\uW_0)\Pi_+=0_4.
    \label{eq:physical-projected-system}
\end{equation}
The first relation is kinematic implying $\uW_0$ cannot contain terms that mix the subspace selected by the factor with its complementary eigenspace. The second applies the Weyl-ordered Liouville equation only within the selected subspace. The off-diagonal projector components are not imposed as additional classical equations.

Using the projector identities above and the fact that the free projectors are independent of $x$, the transport equation becomes
\begin{equation}
    \Pi_+
    \mathcal{L}^{\mathcal{W}}_{\uH_{\mathrm{free}}}(\uW_0)
    \Pi_+
    =
    \frac{p_\sigma^\mu}{m}\PD{\uW_0}{x^\mu}
    =0_4.
    \label{eq:matrix-transport}
\end{equation}
Substitution of Eq.~\eqref{eq:projected-classical-ansatz} therefore gives the full $2\times2$ transport equation
\begin{equation}
    p_\sigma^\mu\PD{\uF_\sigma}{x^\mu}=0_2.
    \label{eq:internal-transport}
\end{equation}
Writing $x^0=ct$ gives
\begin{equation}
    \PD{\uF_\sigma}{t}
    +
    \sigma\frac{c^2p^i}{E_p}\PD{\uF_\sigma}{x^i}
    =0_2,
    \label{eq:coordinate-time-transport}
\end{equation}
which transports arbitrary populations and internal coherences with the correct relativistic branch velocity. It is also a phase-space continuity equation: defining the current-density $\mathbf{J}_\sigma^\mu=(p_\sigma^\mu/m)\uF_\sigma$ gives $\partial_\mu\mathbf{J}_\sigma^\mu=0_2$. The same conservation law holds for the trace and for every fixed matrix component.

The construction is covariant because $S(\Lambda)\Pi_+(p)S^{-1}(\Lambda)=\Pi_+(\Lambda p)$, so both the projected ansatz and Eq.~\eqref{eq:internal-transport} retain their form under a Lorentz transformation. The initial data on a spacelike hypersurface may therefore be any sufficiently regular $2\times2$ matrix distribution $\uF_\sigma$. Hermiticity and positive semidefiniteness may be imposed when $\uF_\sigma$ is interpreted as a statistical density. Equation~\eqref{eq:internal-transport} transports every matrix element without any additional directional constraint. The remaining projector components acquire a different and precise role in the deformation expansion below.

If we impose that the full matrix Liouvillian vanishes, then we find that its off-diagonal blocks impose additional differential constraints on $\uW_0$ that are not required by relativistic transport and restrict otherwise admissible internal populations and coherences~\cite{everitt2026massshellfactorisationspinattempt}.
For this reason we must allow that the full matrix Liouvillian need not vanish which means that it can contain components connecting the subspace selected by the factor constraint with its complementary eigenspace. 
These components will not describe transport within the admissible classical sector, rather, they may be regarded as the matrix analogue of the reaction associated with a classical constraint.
We could therefore write the constrained Liouville equation as
\begin{equation}
    \mathcal{L}^{\uW}_{\uH_{\mathrm{free}}}(\uW_0)
    =\mathcal{R}_0, \text{ where } 
    \mathsf D_{\uH_{\mathrm{free}}}[\mathcal{R}_0]=0_4,
\end{equation}
so $\mathcal{R}_0$ contains only components connecting distinct eigenspaces of $\uH_{\mathrm{free}}$. 
Projecting this equation onto the diagonal blocks eliminates the constraint reaction and, once more, gives
\begin{equation}
    \mathsf D_{\uH_{\mathrm{free}}}
    \left[
   		\mathcal{L}^{\uW}_{\uH_{\mathrm{free}}}(\uW_0)
    \right]=0_4.
\end{equation}
In ordinary constrained mechanics, the reaction required to maintain a constraint is supplied by a Lagrange multiplier. Here the constraint acts in an internal matrix space, and every off-diagonal component connecting distinct eigenspaces of $\uH_{\mathrm{free}}$ can be represented by a commutator with it. 
The corresponding reaction therefore has the form of an internal constraint torque rather than a force in spacetime. 
We will next find that the deformation-quantisation expansion identifies $\mathcal{R}_0=-[\uH_{\mathrm{free}},\uW_1]/\ui$, so this constraint reaction, required to prevent the system from leaving the allowed constraint subspace, is supplied by the first inter-sector correction rather than by an additional classical degree of freedom.

This structure suggests that deformation quantisation may be understood as a hierarchy of constraint-preservation conditions.
For the free generator considered here, the projected part at each order in $\hbar$ supplies the compatibility condition for transport within the allowed sector, while the complementary part determines the inter-sector correction at the next order. 
A complete derivation would require this hierarchy to close consistently for arbitrary phase-space generators, rather than for
$\uH_{\mathrm{free}}$ alone. 
This is therefore left as a subject for future work. 
An outline of such an argument would begin with a formal bidifferential, translation-invariant and associative deformation of ordinary multiplication, requiring a unit, Hermiticity, Weyl symmetry and recovery of the canonical Poisson bracket in the first antisymmetric term. 
The next step would be to verify that the resulting order-by-order consistency conditions lead to the Weyl-Moyal product on flat phase space, at least up to formal equivalence.  
The present free generator cannot by itself determine the full deformation, since its linearity in momentum makes all higher bidifferential terms in its star commutator vanish. 
Consequently, the present calculation does not establish that the Moyal product follows from the constraint hierarchy.
However, deriving it directly from constraint preservation for general generators remains a possible extension of the present construction. 
The calculation below nevertheless shows that the matrix Moyal equation provides a natural generating equation for the hierarchy where its $\hbar$-expansion combines the matrix commutator, which encodes the internal constraint reaction, with the Weyl-ordered Liouvillian, which governs phase-space transport.

Following the Moyal deformation quantisation approach we now deform ordinary multiplication while preserving the order of
the matrix factors~\cite{GROENEWOLD1946405,Moyal_1949}. For matrix-valued phase-space functions $\uF$ and
$\uG$, we define
\begin{equation}
  \uF\star \uG
  =
  \uF\exp\left[
    \frac{i\hbar}{2}
    \left(
      \frac{\overleftarrow{\partial}}{\partial p_\mu}
      \frac{\overrightarrow{\partial}}{\partial x^\mu}
      -
      \frac{\overleftarrow{\partial}}{\partial x^\mu}
      \frac{\overrightarrow{\partial}}{\partial p_\mu}
    \right)
  \right]\uG.
  \label{eq:star-product}
\end{equation}
The zeroth-order term is the ordinary matrix product,
$\lim_{\hbar\rightarrow0}\uF\star \uG=\uF\uG$. The Moyal deformation
therefore introduces both nonlocality in phase space and the
dimensional scale $\hbar$ without changing the underlying Clifford
algebra.


The extension of the Moyal bracket to a matrix star commutator yields
\begin{align}
	  \MB{\mathbf{\uH}}{\mathbf{\uW}} &= \frac{1}{\mathrm{i}\hbar}({\uH} \star {\uW} - \uW \star \uH )\nonumber \\ &= \frac{1}{\ui  \hbar}
      \COM{\mathbf{\uH}}{\mathbf{\uW}} +
	  \mathcal{L}^{\mathcal{W}}_\uH(\uW)+\Order{\hbar}.
\end{align}
For the free generator, which is linear in momentum, this expansion is exact. The natural relativistic matrix Moyal transport equation is $\MB{\uH}{\uW}=0$. To admit a regular classical limit, write
\begin{equation}
    \uW_\hbar
    =
    \uW_0+\hbar\uW_1+\Order{\hbar^2}.
    \label{eq:semiclassical-expansion}
\end{equation}
Substitution into the matrix Moyal equation gives successively
\begin{equation}
    \COM{\uH}{\uW_0}=0_4,
    \label{eq:leading-commutator}
\end{equation}
and
\begin{equation}
    \frac{1}{\ui}\COM{\uH}{\uW_1}
    +
    \mathcal{L}^{\mathcal{W}}_\uH(\uW_0)
    =0_4.
    \label{eq:first-order-hierarchy}
\end{equation}
The leading equation is precisely the classical compatibility condition in Eq.~\eqref{eq:projected-classical-system}. In a basis that diagonalises $\uH$, the diagonal blocks of a commutator with $\uH$ vanish. Applying $\mathsf{D}_{\uH}$ to Eq.~\eqref{eq:first-order-hierarchy} therefore gives
\begin{equation}
    \mathsf{D}_{\uH_{\mathrm{free}}}\!\left[
        \mathcal{L}^{\mathcal{W}}_\uH(\uW_0)
    \right]
    =0_4,
    \label{eq:projected-solvability}
\end{equation}
which recovers the projected classical transport equation. For distinct eigenvalues $\varepsilon_r\ne\varepsilon_s$, the off-diagonal blocks connecting their eigenspaces instead determine
\begin{equation}
    \Pi_r\uW_1\Pi_s
    =
    -\frac{\ui}{\varepsilon_r-\varepsilon_s}
    \Pi_r\mathcal{L}^{\mathcal{W}}_\uH(\uW_0)\Pi_s.
    \label{eq:off-diagonal-correction}
\end{equation}
They are therefore deformation corrections rather than additional differential constraints on $\uW_0$. 

A sufficient, and stronger, condition for the matrix Moyal bracket
to vanish is the pair of left and right stargenvalue equations
\begin{equation}
  \uH\star\uW
  =\epsilon\uW
  =\uW\star\uH.
  \label{eq:left-right-stargenvalue}
\end{equation}
This has the same form as the time-independent Schr\"odinger Wigner/Moyal stargenvalue equations of nonrelativistic quantum mechanics but in phase spacetime. The appearance of $\hbar$ in the deformation now supplies the
physical angular-momentum scale. In the rest frame, the corresponding spin generators
are
\begin{equation}
  S^i
  =\frac{\hbar}{2}\Sigma^i,
  \qquad
  \COM{S^i}{S^j}
  =i\hbar\epsilon^{ijk}S^{\,k}.
  \label{eq:physical-spin}
\end{equation}
The Clifford factorisation therefore supplies the two-state spinor algebra, while deformation quantisation supplies its physical scale.
For $\epsilon=0$, the leading local terms reproduce the projected classical ansatz: $\uH_{\mathrm{free}}\uW_0=0_4=\uW_0\uH_{\mathrm{free}}$. Multiplication by the complementary factor also reproduces the mass-shell constraint. Indeed, from $\uH_{\mathrm{free}}\uW_0=0_4$ we obtain
\begin{equation}
  (\gamma^\mu p_\mu+mc 1_4)
  \uH_{\mathrm{free}}\uW_0
  =
  c(p^\mu p_\mu-m^2c^2)\uW_0=0_4.
  \label{eq:mass-shell-from-star}
\end{equation}
A nonzero leading-order distribution is therefore supported on the mass shell. 
At finite $\hbar$, the full star equations contain gradient corrections and, as we shall soon see, should instead be understood as the
corresponding Dirac--Wigner constraints:
\begin{equation}
\uH \star \uW = 0
= \uW \star \uH. \label{eq:left-right-stargenvalue-epsilon-zero}
\end{equation}
Following~\cite{D1,D2,D3,D4} we seek a direct comparison of the matrix-valued Moyal equation with the Dirac equation by considering the Wigner transform of the Dirac density matrix.  Let $\psi(x)$ satisfy
the minimally coupled Dirac equation
$$
    \left[
        \ui\hbar c\gamma^\mu
        \left(
            \PD{}{x^\mu}
            +
            \frac{iq}{\hbar c}A_\mu
        \right)
        - mc^2
    \right]
    \psi(x)
    =
    0.
$$
The spinor density is then defined as
$
    \rho_{\alpha\beta}(x_1,x_2)
    \equiv
    \psi_\alpha(x_1)\bar\psi_\beta(x_2),
$
where $\bar \psi = \psi^\dag \gamma^0$ is the Dirac adjoint.
We now introduce the centre and relative coordinates
$
\xi=\frac{x_1+x_2}{2}$ and $
y=x_1-x_2.
$
A spinor-matrix Wigner function is then given by performing a Wigner transformation (i.e. a Fourier transform with respect to the coordinate separation) on $\rho_{\alpha\beta}(x_1,x_2)$, namely
$$
W_{\alpha\beta}(\xi,p)
=
\int \ud^4y\,
\exp\!\left(\frac{i}{\hbar}p_\mu y^\mu\right)
\rho_{\alpha\beta}\left(\xi+\frac{y}{2},\xi-\frac{y}{2}\right),
$$
The expression above is the ordinary fixed-gauge Wigner transform.
A fully gauge-covariant Wigner function instead inserts a Wilson line between $x_1$ and $x_2$~\cite{D1,wilsonLine}. Our treatment could therefore be regarded as a fixed-gauge or local-potential version of such treatments.

Applying the Dirac operator to the left coordinate $x_1$ and then
Wigner transforming gives a left star equation
$
H\star W=0,
$
where
$
H(\xi,p)
=
    \gamma^\mu(c p_\mu - qA_\mu(\xi))-mc^2 \, 1_4.
$
Similarly, applying the adjoint Dirac equation to the right coordinate
$x_2$ gives
\mbox{$
W\star H=0
$}.
Thus, we find that the Dirac equation for the spinor amplitude becomes a pair of left and right phase-space equations for the spinor-matrix Wigner
function:
$$
H\star W=0,
\qquad
W\star H=0.
$$
These equations are consistent with the left and right zero-eigenvalue equations in Eq. (\ref{eq:left-right-stargenvalue-epsilon-zero}), with $\epsilon =0$ on the mass shell.
Thus the matrix theory appears to be in agreement with the phase-space Dirac equation. 

To complete the analogy to the previous section, taking the difference of the two Dirac-Wigner equations gives the matrix Moyal-bracket transport equation
$$
H\star W-W\star H=0,
$$
or $\MB{H}{W}=0$, the usual transport equation. Their sum,
$
H\star W+W\star H=0,
$ supplies the companion constraint equation.

We have shown that a first-order linear factorisation of the relativistic mass-shell constraint leads to a Clifford algebra and, in $3+1$ dimensions, to a $4\times4$ matrix-valued phase-space distribution. 
This step is kinematic, requiring neither an operator substitution nor an independent postulate of spin. 
The resulting rank-two on-shell sectors supply the internal two-state structure,
while deformation quantisation introduces the scale $\hbar$ needed to interpret that structure as physical spin-$1/2$.

On each mass-shell sheet, the leading distribution is an arbitrary $2\times2$ matrix within the rank-two kernel and therefore retains internal populations and coherences. Projecting the Weyl-ordered Liouville equation into the same subspace gives the usual relativistic transport equation and its associated continuity law for arbitrary initial matrix distributions. The off-diagonal projector components determine the first deformation correction rather than imposing additional restrictions on the leading classical distribution.
After Moyal deformation, the left and right zero-eigenvalue equations reproduce the constraint and transport structure of the Dirac--Wigner formulation.
The main conclusion is that a first-order relativistic statistical theory naturally carries a spinor structure. Deformation quantisation then provides a phase-space route to the Dirac-Wigner equations and introduces the angular-momentum dimensioned scale $\hbar$.

More speculatively, the present framework suggests that some difficulties in relativistic statistical mechanics may arise from imposing an overly restrictive Hamiltonian structure on a covariant theory. 
A Clifford-valued phase-spacetime formulation may instead provide a setting in which covariance, mass-shell branch structure, and statistical transport coexist, with deformation quantisation
providing a possible route to the corresponding quantum theory.

\section*{Acknowledgements}
I am grateful to those who commented on earlier versions of this work. 
In particular, I am indebted to Tim Spiller and Raymond Bishop for detailed feedback and incisive questions, especially to Tim for his challenge on how to interpret the classical limit of spin in the absence of Planck's constant and Ray for his reviews and suggestions on later drafts of the manuscript. 
I thank John Samson for repeatedly identifying mistakes in early drafts.
I also thank Todd Tilma, Jason Ralph, Russell Rundle, and Kieran Bjergstrom for constructive discussions that helped test the novelty of the ideas and clarify their presentation, and Alexander Balanov and Alexandre Zagoskin for additional stimulating conversations. 
Ben Davies kindly suggested improvements that enhanced the readability of the manuscript.

\section*{CRediT authorship contribution statement}
\textbf{Mark Everitt:} Conceptualisation, Methodology, Formal analysis,
Investigation, Writing -- original draft, Writing -- review \& editing.

\section*{Funding}
This research did not receive any specific grant from funding agencies in the
public, commercial, or not-for-profit sectors.

\section*{Declaration of competing interest}
The author declares that he has no known competing financial interests or
personal relationships that could have appeared to influence the work reported
in this paper.

\section*{Data availability}
No data were generated or analysed for this theoretical study.

\section*{Declaration of generative AI use}
The author acknowledges the use of successive versions of OpenAI's ChatGPT
during the development of this work for checking algebraic consistency,
exploring alternative formulations, identifying possible gaps in reasoning,
and improving the clarity of the presentation. 
The author reviewed and edited
the resulting material and takes full responsibility for the content of the
article.




%\input{appendix_constraint_subspace}

\bibliographystyle{elsarticle-num}
\bibliography{ref}
\end{document}
